%!Mode:: "TeX:UTF-8"
%!TEX TS-program = xelatex
%arara: xelatex
\documentclass{ctexart}
\newif\ifpreface
\prefacetrue
\input{../../../global/all}
\begin{document}
\large
\setlength{\baselineskip}{1.2em}
\ifpreface
\input{../../../global/preface}
\newgeometry{left=2cm,right=2cm,top=2cm,bottom=2cm}
\else
\newgeometry{left=2cm,right=2cm,top=2cm,bottom=2cm}
\maketitle
\fi
%from_here_to_type
如无特殊说明，我们总认为
\[
  X=\{\Omega,\mathscr{F},\mathscr{F}_n,X_n,\theta_n,\mathbb P^x\}
\]
是时间集 \(T=\mathbb N\)、状态空间 \(E\) 离散的 Markov 链。其转移函数为
\(p(x,y)\)。

对于任意 \(x\in E\)，记
\[
  \tau_x^+ := \inf\{n\ge 1:X_n=x\}
\]
表示 \(x\) 的首中时。定义
\[
  \tau_x^{(1)}:=\tau_x^+,
  \qquad
  \tau_x^{(k)}
  :=
  \inf\{n>\tau_x^{(k-1)}:X_n=x\},
  \quad k\ge2.
\]

\begin{problem}
  证明：\(\tau_x^{(k)}\) 是 \(\mathscr F_n\)-停时。
\end{problem}

\begin{solution}

  我们对 \(k\) 做归纳。

  首先，
  \[
    \tau_x^{(1)}
    =
    \tau_x^+
    =
    \inf\{n\ge1:X_n=x\}.
  \]
  对任意 \(n\ge1\)，有
  \[
    \{\tau_x^{(1)}=n\}
    =
    \{X_1\neq x,\dots,X_{n-1}\neq x,X_n=x\}.
  \]
  由于 \(X_i\) 是 \(\mathscr F_i\)-可测且
  \(\mathscr F_i\subseteq\mathscr F_n\)，故
  \[
    \{\tau_x^{(1)}=n\}\in\mathscr F_n.
  \]
  于是
  \[
    \{\tau_x^{(1)}\le n\}
    =
    \bigcup_{m=1}^n
    \{\tau_x^{(1)}=m\}
    \in\mathscr F_n,
  \]
  因此 \(\tau_x^{(1)}\) 是停时。

  假设 \(\tau_x^{(k-1)}\) 是停时。考虑
  \[
    \tau_x^{(k)}
    =
    \inf\{n>\tau_x^{(k-1)}:X_n=x\}.
  \]

  对任意 \(n\)，有
  \[
    \{\tau_x^{(k)}=n\}
    =
    \bigcup_{m=1}^{n-1}
    \Big(
      \{\tau_x^{(k-1)}=m\}
      \cap
      \{X_{m+1}\neq x,\dots,X_{n-1}\neq x,X_n=x\}
    \Big).
  \]
  由归纳假设，
  \(\{\tau_x^{(k-1)}=m\}\in\mathscr F_m\subseteq\mathscr F_n\)，
  而后面的事件显然属于 \(\mathscr F_n\)，故
  \[
    \{\tau_x^{(k)}=n\}\in\mathscr F_n.
  \]
  于是
  \[
    \{\tau_x^{(k)}\le n\}
    =
    \bigcup_{r=1}^n
    \{\tau_x^{(k)}=r\}
    \in\mathscr F_n.
  \]
  因此 \(\tau_x^{(k)}\) 是 \(\mathscr F_n\)-停时。

\end{solution}

\begin{problem}\label{pro:1}
  假设 \(x\) 可达 \(y\)，且 \(x\) 常返。证明：

  \begin{enumerate}[label=(\arabic*)]
    \item
      \[
        \mathbb P^x(\tau_y^+<\infty)
        =
        \mathbb P^y(\tau_x^+<\infty)
        =
        1.
      \]
      特别地，\(x,y\) 互达。

    \item
      \(y\) 也常返。
  \end{enumerate}
\end{problem}

\begin{solution}

  由于 \(x\to y\)，存在某个 \(m\ge1\)，使得
  \[
    p^{(m)}(x,y)>0.
  \]

  \begin{enumerate}
    \item
      因为 \(x\) 常返，
      \[
        \mathbb P^x(\tau_x^+<\infty)=1.
      \]
      并且从 \(x\) 出发会无穷次返回 \(x\)。

      每次回到 \(x\) 后，在接下来的 \(m\) 步内到达 \(y\) 的概率至少为
      \[
        p^{(m)}(x,y)>0.
      \]
      利用强 Markov 性，链每次回到 \(x\) 后都会重新开始，因此若始终不访问 \(y\)，其概率至多为
      \[
        \lim_{n\to\infty}(1-p^{(m)}(x,y))^n=0.
      \]
      故
      \[
        \mathbb P^x(\tau_y^+<\infty)=1.
      \]

      接下来证明
      \[
        \mathbb P^y(\tau_x^+<\infty)=1.
      \]

      设
      \[
        a=\mathbb P^y(\tau_x^+<\infty).
      \]
      由全概率公式，
      \[
        1
        =
        \mathbb P^x(\tau_y^+<\infty)
        \le
        \mathbb P^x(\tau_y^+<\infty,\tau_x^+\circ\theta_{\tau_y^+}<\infty).
      \]
      利用强 Markov 性，
      \[
        \mathbb P^x(\tau_x^+\circ\theta_{\tau_y^+}<\infty
          \mid
        \mathscr F_{\tau_y^+})
        =
        \mathbb P^y(\tau_x^+<\infty)
        =a.
      \]
      因此
      \[
        1\le a,
      \]
      故 \(a=1\)。

      于是
      \[
        \mathbb P^y(\tau_x^+<\infty)=1.
      \]
      所以 \(y\to x\)，即 \(x,y\) 互达。
    \item
      由于已经证明 \(y\leftrightarrow x\)，故存在 \(n\ge1\) 使得
      \[
        p^{(n)}(y,x)>0.
      \]

      从 \(y\) 出发，几乎必然到达 \(x\)，而 \(x\) 出发几乎必然回到\(y \)，
      故：
      \[
        \mathbb{P}^y(\tau_y^+<\infty)\geq \mathbb{P}^y(\tau_x^+<\infty, \tau_y^+ \circ \theta_{\tau_x^+}<\infty ) = \mathbb{P}^y(\tau_x^+<\infty) \mathbb{P}^x(\tau_y^+<\infty) = 1
      \]。

      故
      \[
        \mathbb P^y(\tau_y^+<\infty)=1,
      \]
      即 \(y\) 常返。
  \end{enumerate}

\end{solution}

\begin{problem}
  假设 \(x\) 常返。证明：对于任意 \(k\ge1\),
  \[
    \mathbb P^x(\tau_x^+\circ\theta_k<\infty)=1.
  \]
\end{problem}

\begin{solution}

  由于 \(x\) 常返，
  \[
    \mathbb P^x(\tau_x^+<\infty)=1.
  \]

  考虑时刻 \(k\) 之后的过程。由 Markov 性，
  在给定 \(X_k=y\) 的条件下，
  \[
    \mathbb P^x(\tau_x^+\circ\theta_k<\infty\mid X_k=y)
    =
    \mathbb P^y(\tau_x^+<\infty).
  \]

  由于 \(x\) 常返，由\ref{pro:1}可得
  \[
    \mathbb P^y(\tau_x^+<\infty)=1
  \]
  对所有满足 \(\mathbb{P} (y=X_k)>0\) 的\(y \) 成立。

  于是
  \[
    \mathbb P^x(\tau_x^+\circ\theta_k<\infty)
    =
    \sum_{y\in E}
    \mathbb P^x(X_k=y)
    \mathbb P^y(\tau_x^+<\infty)
    =
    \sum_{y\in E}\mathbb P^x(X_k=y)
    =
    1.
  \]

  故结论成立。

\end{solution}
\end{document}
